
\section{The Design Trade-off Matrix}
\label{sec:tradeoff}

\Cref{tab:tradeoff} consolidates the chapter. Each row is an option,
each option is scored on the four properties that the theory of
\Cref{ch:foundations} identifies as the relevant dimensions, and the
final column names the behaviour the option induces. The scoring is
ordinal and theoretical, not measured; it summarises the arguments made
above rather than adding to them.

\begin{sidewaystable}
\centering
\caption[The FBA design trade-off matrix]{The FBA design trade-off
matrix. Ratings are ordinal and derived from theory:
$+$ favourable, $\circ$ neutral or conditional, $-$ unfavourable.}
\label{tab:tradeoff}
\small
\begin{tabularx}{\textwidth}{@{}p{2.2cm} p{3.4cm} c c c c L@{}}
\toprule
\textbf{Block} & \textbf{Option} & \textbf{Speed neutral.} &
\textbf{Price disc.} & \textbf{Exec.\ certainty} & \textbf{Impl.\
simplicity} & \textbf{Induced behaviour}\\
\midrule
\multirow{4}{2.2cm}{Interval}
 & Short $\tau$ ($\leq 10$\,ms)      & $\circ$ & $+$    & $+$    & $\circ$ & Thin batches; residual speed value\\
 & Long $\tau$ ($\geq 500$\,ms)      & $+$     & $-$    & $-$    & $+$     & Stale orders; delay cost\\
 & Fixed boundary                    & $\circ$ & $+$    & $+$    & $+$     & Boundary sniping\\
 & Randomised close                  & $+$     & $\circ$& $-$    & $-$     & No timing target; planning cost\\
\midrule
\multirow{4}{2.2cm}{Disclosure}
 & Fully open book                   & $-$     & $+$    & $+$    & $+$     & Intent exposure; late submission\\
 & Indicative price                  & $\circ$ & $+$    & $\circ$& $\circ$ & Coordination without order-level leakage\\
 & Aggregate imbalance               & $+$     & $\circ$& $\circ$& $\circ$ & Directional signal only\\
 & Sealed bid                        & $+$     & $-$    & $-$    & $+$     & Bid shading\\
\midrule
\multirow{3}{2.2cm}{Segregation}
 & Unified pool                      & $\circ$ & $+$    & $+$    & $+$     & Cross-subsidy of informed flow\\
 & Class segregation                 & $\circ$ & $\circ$& $\circ$& $-$     & Tighter quotes to benign flow; label gaming\\
 & Intent segregation (DFBA)         & $+$     & $\circ$& $-$    & $-$     & No maker-vs-maker trades; split liquidity\\
\midrule
\multirow{5}{2.2cm}{Rationing}
 & Price-time in batch               & $-$     & $\circ$& $+$    & $+$     & Micro speed race\\
 & Pro rata                          & $+$     & $\circ$& $\circ$& $+$     & Order size inflation\\
 & Size priority                     & $+$     & $\circ$& $\circ$& $+$     & Threshold clustering\\
 & Random                            & $+$     & $\circ$& $-$    & $\circ$ & Unbiased size; execution uncertainty\\
 & Hybrid $\kappa$                   & $\circ$ & $\circ$& $\circ$& $-$     & Attenuated distortions\\
\midrule
\multirow{3}{2.2cm}{Settlement}
 & Gross transfers                   & $\circ$ & $\circ$& $+$    & $+$     & Full per-fill auditability\\
 & Net settlement                    & $\circ$ & $\circ$& $+$    & $-$     & Lower transfer cost; opaque per-fill mapping\\
 & Roll-over residual                & $\circ$ & $+$    & $\circ$& $+$     & Standing liquidity accumulates across batches\\
\bottomrule
\end{tabularx}
\end{sidewaystable}

\section{Case Studies: Deployed Batch Auction Systems}
\label{sec:casestudies}

The taxonomy is worth testing against systems that exist.
\Cref{tab:casestudies} places five of them in the design space.

\textbf{Equity opening and closing auctions} are the oldest deployment
and the most conservative point in the space: a long interval, an
indicative price, a unified pool, and price-time rationing, with a
randomised extension to protect the boundary. They demonstrate that
uniform-price clearing is operationally routine at scale, but their
frequency is once or twice per day, so they neutralise nothing about the
intraday speed race.

\textbf{The Taiwan Stock Exchange} operated intraday call auctions at
short intervals for decades before moving to continuous trading, which
makes it the closest thing to a natural experiment on frequent batching
in equities \citep{lee2026taiwan}.

\textbf{Jump Trading's DFBA} is the clearest instance of intent-based
segregation combined with batching, discussed in
\Cref{sec:segregation}.

\textbf{CoW Protocol} clears batches of user intents with a uniform
clearing price per token pair, delegating the clearing computation to a
competitive solver auction \citep{dutta2025cow,harrison2025hybrid}. It
is the deployed system in which the multi-asset formulation of
\eqref{eq:multiflow} does the most work, since coincidence-of-wants
matching within a batch avoids touching external liquidity at all.

\textbf{SPEEDEX} demonstrates that arbitrage-free multi-asset batch
clearing can be computed fast enough for a high-throughput venue, which
addresses the standard objection that batch clearing is too slow to be
practical \citep{ramsay2023speedex}.

\begin{table}[htbp]
\centering
\caption[Deployed batch auction systems in the design space]{Five
deployed systems placed in the design space of \Cref{sec:blocks}.}
\label{tab:casestudies}
\small
\begin{tabularx}{\textwidth}{@{}p{2.6cm} p{2.0cm} p{2.4cm} p{2.3cm} L@{}}
\toprule
\textbf{System} & \textbf{Interval} & \textbf{Disclosure} &
\textbf{Segregation} & \textbf{Clearing and rationing}\\
\midrule
Equity open/close auctions & Once or twice daily, randomised extension &
Indicative price and imbalance & Unified & Volume maximisation,
price-time rationing\\
TWSE intraday call auctions & Seconds to tens of seconds & Partial &
Unified & Volume maximisation\\
Jump DFBA & Sub-second batches & Not fully public & Maker/taker intent
partition & Two auctions per batch, one uniform price\\
CoW Protocol & Block-paced batches & Public intents & Unified, solver
mediated & Multi-asset uniform clearing via solver competition\\
SPEEDEX & Per ledger/block & Public & Unified & Multi-asset
arbitrage-free uniform clearing\\
\bottomrule
\end{tabularx}
\end{table}

\section{Answer to RQ1}
\label{sec:rq1answer}

RQ1 asked how a frequent batch auction engine is designed and which
options exist for each of its architectural components. The answer
developed in this chapter has three parts.

An FBA engine is fully specified by five design blocks: the batch
interval, the disclosure regime, the segregation architecture, the
clearing and rationing rule, and settlement and residual handling.
\Cref{fig:pipeline} shows that these blocks are sequential stages of one
batch cycle, and that the corresponding decisions in a continuous engine
are not absent but implied by serial processing, which is why they
become visible only once batching is introduced.

Within each block the options are enumerable and their behavioural
consequences derivable. \Cref{tab:tradeoff} states them. The recurring
structure is that each option shifts competition from one variable to
another rather than eliminating it: randomising the close removes the
timing target but adds planning cost, opacity removes intent exposure
but adds bid shading, pro rata removes the timing channel but adds size
inflation. The only design choice in this space that is close to a free
improvement is uniform pricing itself, and even that has a cost, namely
that the marginal participants at $\clr$ must be rationed by some rule
that will in turn be gamed.

Finally, two questions in this space have no settled answer and are
contributions of this chapter rather than summaries of the literature.
The maker--taker distinction has no natural definition under
simultaneous uniform-price clearing, and four candidate redefinitions
were set out in \Cref{sec:segregation}, each of which subsidises
different behaviour. And the price improvement that uniform pricing
produces relative to submitted limits is not a fixed property of the
mechanism but a quantity whose size and distribution across participants
depends on the clearing and rationing rules chosen, which is precisely
what the empirical part of this thesis measures.
